\scanpage{30}
\uline{Dém} En deux étapes

\uline{\unsure{Lemme|Lem}, 1}, $\xi_k=0$ si \inserted{$k\geq N$ tel que} $c_{i,k-i}\not\equiv0\ (p)$.

\uline{Lemme} 1.1. Soient $\deleted{\unsure{k',k|k_1,k}}\ k'\geq k\geq N$ tels que, posant $j=k-i$, $j'=k'-i\deleted{\unsure{\prime}}$, on ait
\[
c_{i,j}\not\equiv0\ (p)\quad\text{et}\quad c_{i,j'}\not\equiv0\ (p)\tag*{(4)}
\]
Alors on a l'équivalence
\[
c_{j'-j,j}\not\equiv0\ (p)\iff c_{j'-j,j+\deleted{\unsure{j}}i}\not\equiv0\ (p)\tag*{(5)}
\]
et si les conditions équivalentes (5) sont satisfaites, on a l'équivalence
\[
\xi_k=0\iff\xi_{k'}=0.\tag*{(6)}
\]

\uline{Dém de} 1.1. Calculons dans $\mathbb Q$
\[
\begin{aligned}
\text{(a)}\quad &c_{j'-j,j+i}\,/\,c_{j'-j,j}=\frac{(j'+i)!}{\deleted{(j'-j)!}(j+i)!}:\frac{(j'\deleted{\unsure{+j-j}})!}{\deleted{(j'-j)!}j!}\\
&\text{et }\deleted{\unsure{/}}\\
\text{(b)}\quad &c_{i,j'}\,/\,c_{ij}=\frac{(i+j')!}{\deleted{i!}j'!}:\frac{(i+j)!}{\deleted{i!}j!}
\end{aligned}
\]
\deleted{\unsure{et}}
\deleted{\[
\text{(b)}\quad c_{i,j'}/c_{ij}=\frac{(i+j')!}{i!\,j'!}:\frac{(i+j)!}{i!\,\unsure{j!}}
\]}
\deleted{Or dans $\mathbb Z_p$, pour la \unsure{rel.|val.} de divisibilité, on \unsure{le quotient de (a) et (b) est} grâce à (4),
\[
(i+j)!\sim i!\,j!\qquad (i+j')!\simeq i!\,j'!
\]
donc le premier \unsure{membre} de (a) est \unsure{eq} à
\[
\frac{i!\,j'!}{i!\,j!}:\frac{(j'+i-j)!}{i!}=
\]}
Donc les deux \unsure{quantités|quotients} (a) et (b) sont égaux, et comme (b) est un quotient de deux unités \deleted{$p$-adiques} \inserted{de $\mathbb Z_p$}, on voit que $c_{j'-j,j+i}$ et $c_{j'-j,j}$ sont équivalents pour la divisibilité, donc l'un est une unité sss l'autre l'est. Quand il en est ainsi, la
% archive sheet number: 29

